\section{Limitations, external validity, licensing, and escalation}
\label{sec:limitations}

\emph{Status.} Most of this section scopes the deferred real-agent program, and
is retained so its limitations are on record; the completed evidence is the
closed-fixture and synthetic work of \cref{sec:main-measurements,sec:main-primary}.
The live substrate below has not been qualified and no confirmatory result
exists.  The closed-fixture and synthetic evidence is external-validity narrow in
the same directions stated here, and the negative cost result is not a
cost-superiority claim.

The intervention is one-step and all values depend on the declared continuation
policy, horizon, agent/model version, fault mechanism, task frame, availability
rules, and adjudication procedure.  Two scope boundaries follow from the
sampling scheme rather than from the substrate, and are stated separately
because they survive a perfect instrument.  First, \emph{deployment timing}: the
checkpoint is drawn uniformly from a completed trajectory's realised eligible
window, so \cref{prop:l0} certifies a controller that intervenes once at a
checkpoint drawn that way.  That is implementable online only when the window is
known in advance; a stopping-time rule --- ``intervene at the first checkpoint
whose cell is positive'' --- adds a selection term the design does not estimate,
so any online reading of the guarantee is scoped to restored-checkpoint
intervention.  Second, \emph{successor coverage}: the design restores the
sampled checkpoint and never a state reached after an intervention, so the span
of the one-step advantage over post-intervention states is decision-relevant and
unmeasured (\cref{prop:transfer,rem:transfer-fails}).  A depth-two pilot would
sample a \emph{declared} successor distribution; it would not establish a global
advantage span over all reachable states, and no such bound is claimed.  Neither this design nor predictive accuracy
establishes representation sufficiency; \cref{thm:refinement} bounds the
resulting loss from below but cannot certify its absence.  Full-arm branching
estimates regime means for the sampled stochastic execution process.  It does
not identify individual treatment effects, the fraction of individual
\emph{trajectories} harmed, or the benefits of repeated adaptive intervention.
It does identify the fraction of \emph{checkpoints} at which a regime is harmful
in expectation, \eqref{eq:harm-fraction}, because every arm is executed at the
same checkpoint; the sharp non-identification results for the harmed fraction
concern designs that assign one arm per unit and do not bind here.  Unequal or
selective regime availability can restrict comparisons to small common domains.

The one-step restriction also bounds interaction cost: production orchestration
work reports that online exploration and training/evaluation budgets can be
hard to bound under strict API-cost constraints \citep{gstar2026}.

The theoretical results are individually elementary and are presented as such.
\Cref{thm:safe} is a one-step instance of a pessimistic safe-improvement
argument, \cref{prop:certification} a union bound, \cref{thm:ordinal} a
rearrangement inequality.  What is claimed is the object they are applied to and
the observation that the one-step design makes their hypotheses verifiable
rather than assumed.  A reader who expects new machinery will not find it.
One verifiability cost is stated rather than hidden: the safe-improvement gate
operates on a registered finite partition of the representation space
(\cref{eq:gate-family}), so its guarantee is uniform within each cell, and the
cell count $B$ trades gating precision against the multiplicity it adds to the
sample ceiling.

The signed initial substrate and any approved replacement narrow the empirical
population to autonomous ML experimentation.  Concretely, the confirmatory tier
is a single synthetic fault family (F2-FEHPO) on one version-bound agent harness
against a single bound inference substrate, with MLGym excluded and F1-ORD and F2-ADV retained only
as historical contrasts; the title asks about autonomous experiment agents in
general, but the deferred confirmatory evidence would bear on this one
substrate only.  It cannot support claims about other sciences,
signal analysis, or deployment systems without prospective adapters and new
validation.  Public tasks also create contamination risk, which requires
version pinning, disclosure of model training-date limitations where known, and
avoidance of protected evaluation material.  Benchmark code and each dataset
must be used under their own licenses.  MLGym is excluded from the current
protocol; its noncommercial restriction remains a material constraint on any
future interface-transfer study rather than a footnote.

Model, quantization, hardware, decoding, prompt, tools, and harness are parts of
the instrument.  History replay cannot restore undeclared process or cache
state; E0 probes this risk, but its verdict is version-bound.

How strong ``version-bound'' can be is itself substrate-dependent, and this is
why the study serves its model locally.  Under local inference the instrument is
byte-identical for the life of the study and every layer named above is bound by
digest.  A hosted endpoint admits no such binding: the provider can change
weights, quantization, serving stack, or routing behind a stable model
identifier, and behavioural fingerprinting from a fixed prompt set detects that
change after the fact rather than preventing it, with real
provider-to-provider and within-provider variation reported for identical
nominal models \citep{fingerprints2026}.  Under a hosted substrate the
reproducibility claim would be narrower --- a change detected mid-study
invalidates the affected stage rather than being repairable --- and the inference
cache would stop being a layer the evaluator can disable or isolate.  We
evaluated such a substrate and rejected it on those grounds, at the cost of a
smaller model than a hosted endpoint would have supplied.  That trade is the
substantive limitation of the choice: control and auditability were bought with
capability, and the population narrows accordingly.

Local execution does not imply deterministic generation, and no such claim is
made.  Reduction order in accelerator kernels, batching, sampler state, and
server-side prefix caching can each change a completion under a fixed prompt,
seed, and decoding configuration; model sampling randomness is one of the two
declared non-restorable layers for exactly this reason.  The protocol therefore
registers a repeated-request test --- one fixed prompt, seed, sampling
configuration, and tool result replayed across cold servers, a warm server with
reuse disabled, and concurrent workers --- and reports the achieved exact-match
rate rather than asserting reproducibility.  That rate is also the floor against
which E0-M is read: two correctly isolated siblings may differ for reasons
unrelated to leakage, so a sibling-divergence signal is interpretable only
relative to the substrate's own reproducibility.

One structural property should be stated apart from any claim about results.  A
branch is executed by one worker against one model instance, so accelerator
count is a throughput parameter: adding independently isolated workers can
shorten elapsed time without changing the estimand, allocation rule, or
analysis.  That does not extend to accelerator \emph{class}.  A wall rail
converts generation speed into a condition on measurability, so a candidacy
verdict under a fixed elapsed-time contract is jointly a property of the model
and the device, and model, quantization, runtime, harness and accelerator class
are all bound together.

The completed candidacy exam (\cref{app:model-qualification}) ran twelve
candidates on two leased 48~GiB RTX 6000 Ada GPUs and selected Devstral Small
2 24B \emph{provisionally}; earlier development rounds are archived in the evidence record and are not
pooled with it.  Qwen3-8B
remains the recorded engineering control until a successor decision replaces
it.  A provisional candidacy selection is not a substrate qualification:
nothing in that exam measured restoration, isolation, determinism or resource
ceilings, and E0 cannot open until the model smoke, the realistic soak, the
derived ceilings, the replacement R5-2A signature and a separate execution
authorization are all complete.  G1 is unchanged --- at least 200 valid
completed branch triples, restoration rate at least 0.95, and zero detected
leakage.

Three limits of the exam's own design survive it.  Thirty roots resolve a
success rate only to about three roots, and each root is adjudicated on 40
hidden rows, so the screen decides eligibility and licenses no ranking of
models (\cref{app:model-qualification}).  Requiring a model to reach
75\% of a registered reference procedure before any restoration measurement may be
attempted conflates task competence with instrument observability, and
separating operational E0 eligibility from task-yield ranking remains an open
design item.  The tabular family is also the easy case: the workloads that
would test the instrument where it matters --- multi-file coding tasks and
notebook-style analyses, against an agentic screen such as ASI-Bench --- are
named here as the planned harder material and are not part of this round's
registration.

The accessibility statement is correspondingly narrow.  The complete local
instrument is controllable on a two-GPU laboratory allocation, and independent
workers make accelerator count primarily a throughput parameter.  A separate
development program also exercises the harness on four laboratory-owned legacy
12~GiB accelerators, one worker per card; that demonstrates feasibility on the
hardware the authors have, not reproducibility across small laboratories, and it
is not the confirmatory substrate.  Accelerator count is a throughput parameter
only where each branch is served by one model instance: a two-card layer-split
server makes two cards jointly implement one branch, which is a different
deployment rather than a faster one.  Total GPU-hours, wall time, peak memory,
storage, software and licence requirements will be reported from the soak and
E0 rather than projected from candidacy alone.

As of 2026-09-08 the realistic soak has been \emph{executed} on the local
two-GPU deployment and the E0 run has not.  The soak is engineering
qualification of the instrument's substrate, is registered as non-qualifying,
and produces no scientific evidence: 256 synthetic roots and 768 branches, all
settled, with zero prompt-cache reuse measured over $6.56\times10^{7}$ prompt
tokens, a peak of 21{,}135~MiB on one card, 172.9~MB of workspace storage and
3.098~kWh.  Two facts about it qualify every cost figure quoted from this
round.  First, accelerator time is not one quantity: the artifact's
\texttt{gpu\_hours} is summed client-side branch latency across concurrent
workers (14.512 worker-hours), whereas utilisation-weighted card time over the
same interval is 9.138 card-hours and lease occupancy is 16.642 card-hours;
these are three different measurements and are not interchangeable.  Second,
the run exceeded its own frozen accelerator cap on two of those three readings,
which is recorded as a protocol deviation rather than resolved by relabelling
the cap.  The measured slots derived from it remain \emph{unbound} pending
author decisions, and the serving runtime has since been repaired, so the
deployment is not described here as qualified.  E0 remains not run and no
pilot or confirmatory result exists.

Responsible escalation is an intended safe terminal response when a regime's
registered rule concludes that continuation cannot be justified.  It is not
scientific success and does not receive primary completion credit, but it must
never be merged with producing an invalid scientific output.  Rates, latency,
human burden, and downstream disposition of escalation are reported separately.
The final terminal-utility ordering remains an author decision before
confirmatory registration.

\paragraph{Future work.}
Three extensions follow from the contract rather than from any result.  Given
cheap continuation labels and expensive paired labels, how a fixed branch budget
should be split is the same two-phase sampling problem
\cref{sec:main-measurements} already declares open
\citep{neyman1938,angelopoulos2023ppi}; it is stated as an extension of that
problem, not a new one.  Whether a model trained on corrupted replay labels
inherits a systematic allocation error needs a defective pool far larger than the
E0-M extension carries.  Richer representations --- a text encoder over $H$, and
transfer across templates, fault families and harnesses --- are out of scope: one
synthetic fault family on one version-bound substrate identifies neither.  A
transfer claim would also have to say which property is claimed to transfer:
ranking quality and calibration are separate axes and move independently under
distribution shift \citep{damani2025rlcr}, so a monitor that orders checkpoints
well on an unseen fault family may still be badly calibrated on it, and the
converse.  The positioning against rollout-supervised value heads and process
reward models \citep{mathshepherd2024,agentprm_choudhury2025,agentprm_xi2025},
which the draft already distinguishes as value targets under a completer policy,
is left to a fuller treatment.  No results are promised.
